\section{Introduction}
\label{sec:introduction}

Post-compromise recovery depends on the secrecy of contributions incorporated
into later keys. Ratcheting protocols make this dependence precise through the
evolution of key-exchange and messaging state
\cite{signal2020,tripleratchet2025}. In a persistent implementation, a
contribution also has a storage history: its randomness is sampled, a replay
record is saved, the cryptographic operation is executed, and its result is
published and confirmed. These events can be separated by a crash or a delayed
message. A security experiment must account for the secret source across this
entire history.

\paragraph{The reservation boundary.}
Consider an endpoint that durably reserves the random tape of a KEM operation
so that an interrupted call can resume with its original inputs. An adversary
acquires this reservation, its wrapping key, and the predecessor root before
the key-generation call executes. After a crash, the endpoint resumes the call
and completes an honestly authenticated exchange. The captured tape determines
the same private key; the public response then determines the KEM contribution
and, with the predecessor root, the descendant traffic keys. The resumed
computation is a later event in the lifetime of the original secret source.

Future-key exposure already has a precise place in messaging security.
SCKA's vulnerable message sets include retained future-epoch keys
\cite{scka2025}, and Triple Ratchet accounts for next-epoch material sampled
while its transmission is in progress \cite{tripleratchet2025}. The question
here concerns the persistent objects from which such material can be recovered:
\emph{which facts about an original operation and the acquisition of its records
must a recovery experiment preserve through restart, copying, and later key
disclosure?}

There are two related difficulties. First, one reservation may occur in a
pending target, a current journal image, and several retained copies. A later
wrapping-key disclosure can make an earlier capture readable, even after the
honest endpoint has deleted its local copy. Second, a cryptographic reduction
may possess a selected operation's public outputs without possessing its
random tape. It must still account for public record commitments, complete-record
equality, exact retries, reached persistence states, and authenticated reports
whose inputs depend on retained private values. These are semantic and
simulation obligations of the recovery argument.

\paragraph{Original sources and durable executions.}
We associate each reservation with its original owner, operation, session,
epoch, and sampling event. We separately identify an immutable session-record
version and a physical image carrying that version. Copying or resealing an
image preserves the referenced source and record version. Acquisition records
the objects obtained by the adversary; readability is evaluated against the
complete acquisition and key-disclosure history. This representation makes
the secrecy premise stable under changes to physical storage representation
while allowing adversarial knowledge to grow.

We study this representation in Q-Periapt Continuity, a hybrid post-quantum
messaging construction with persistent rekey operations and asynchronous
confirmation. The construction provides a concrete setting in which the
reservation, execution, and confirmation events have distinct operational
boundaries. Its journal can also aggregate several sessions, making the
distinction between a new image and an existing source directly observable.

\paragraph{Main result.}
Theorem~\ref{thm:native-online-message} gives confidentiality and fresh-message
integrity for an online selected application send after the next rekey of an
established session, under its finite interfaces and primitive assumptions.
Each endpoint has one authoritative history with passive copies; the network
and recovery schedule remain adversarial. The predecessor
root and wrapping keys may be known. Earlier message keys and the declared
successor values may be disclosed. The target step is selected before its
first private computation at either endpoint, and the terminal selector may
choose the recovery cut after interaction. Eligibility requires both original
KEM reservations to follow that cut, the target signing secret keys and
designated target-epoch entropy to remain undisclosed, and every acquired
target-session version to have been first saved at or before the cut
(\eqref{eq:eligibility}). The theorem bounds the unconditional eligible-message
advantage and fresh-installation event by explicit source-attribution, record,
KEM, shared-oracle key-schedule and AEAD terms. Its executed-slot selector
connects attribution to ordinary hedged ML-DSA security. The first-message and
selected-root confirmation results are established along the same proof route
(Corollary~\ref{cor:live-core-next-rekey} and
Theorem~\ref{thm:native-selected-authentication}).

\paragraph{Technical contributions.}
The construction connects an original-source exposure condition to this
native message experiment through simulatable private retained state.
Reservations keep their original operations, and acquisition is recorded
separately from later readability. Alias normalization preserves the public
view and original-record ledger under copying and resealing
(Lemma~\ref{lem:alias-normalization}). For fixed public key derivations, a
guarded-path characterization exactly matches iterative decryption
(Theorem~\ref{thm:guarded-normalization}). Conjunctive guards capture how
acquired ciphertexts and disclosed keys jointly reveal a source, while
alternative copies provide additional routes to the same source.

The central reduction constructs retained private state around an operation
from its public outputs, while its random tape remains hidden from the simulator.
Images and intents preserve original record identities; dependent immutable
reports are simulated in reverse causal order. Candidate-local accounting
localizes the original native score before comparing each source pair,
using provenance to justify permanent stops for every eventual recovery cut.
The original-operation adapter carries the resulting execution through the
complete ContextBound contribution, root initialization, confirmation and
traffic schedule.

Online message selection requires an additional dependency argument.
Successive Extract and Expand comparisons isolate the selected key while
preserving earlier-key observations, its disclosed successor seed and real
suffix computation. Each comparison retains the first private computation
across endpoint aliases, including failed receive-side advances. The
composition preserves the accepted send's original sources and candidate-local
record stops with explicit lifetime and whole-execution budgets.
Section~\ref{sec:computation} summarizes these allocation boundaries and the
connections supplying the main result.

Component bridges connect operational ML-KEM to ordinary IND-CCA security
through a deferred-challenge comparison, and reserved ML-DSA randomizers to
ordinary hedged unforgeability through a shared-XOF comparison
(Proposition~\ref{prop:kem-deferred-challenge} and
Theorem~\ref{thm:native-mldsa-reservation-bridge}). Joint entry comparisons
retain the multiple uses of a contribution. Complementary bootstrap instances
protect the original inventory source or reply source and establish historical
activation of the complete signed transcript
(Corollary~\ref{cor:bootstrap-initial-terminal} and
Theorem~\ref{thm:bootstrap-terminal}). A further purpose-key instance retains
equal-input aliases and separate returned-owner lifecycles
(Section~\ref{sec:sdk-purpose-connection}).

A mechanized symbolic theorem uses the same provenance condition for message
secrecy across unbounded epochs (Theorem~\ref{thm:symbolic}). Its model includes
inherited roots, shared key identities, alternating proposers and asynchronous
confirmation. Native executions exhibit the exposure distinctions: traffic is
reconstructed from reservations captured before restart; a delayed-receipt
execution exposes epoch-two traffic while the peers' confirmed counters remain
zero and one; and a sibling-session write carries an earlier target reservation
into a new aggregate image.

\paragraph{Supporting validation.}
Independent finite checks validate the guarded characterization under
acquisition, deletion and late disclosure. Source correspondence binds the
native operations to the declared interfaces. Pinned SPQR and whole-KEM
corpora measure reconstructible message sets across schedules. A further 49
native/reference executions share directed transmission opportunities,
charging control traffic and exact retries to the same budget and recording
reconstructed messages, commitments, deliveries, bytes and persistence work.
A one-way control records exhausted-budget suspension. In a paired 28-run
native ablation, a minimum retry age of four public sender turns increases
completed service in every measured bidirectional or resumed schedule while
retaining the alternating current-epoch exposure counts and recording
schedule-dependent communication costs.

\paragraph{Security model and organization.}
Crash recovery resumes retained, bound operations under the authentication,
acquisition and disclosure interface of Section~\ref{sec:security-model}.
The symbolic theorem covers unbounded epochs. The computational results
connect the stated finite record, report and online message interfaces within
the common message experiment, retaining their explicit cryptographic and
disclosure terms. Section~\ref{sec:construction} describes the construction;
Sections~\ref{sec:security-model} and~\ref{sec:computation} develop the record
model and computational connection. Section~\ref{sec:established} presents
the mechanized recovery theorem and separating executions, and
Section~\ref{sec:evaluation} validates the exposure semantics and presents
the matched executions. Section~\ref{sec:discussion} develops their
consequences and further instantiation.
